\section{Datasets, preprocessing and experimental protocol}
\label{sec:data}
\textbf{FER2013 provenance and protocol:} FER2013 was introduced by \textcite{goodfellow_challenges} as part of a Kaggle challenge. Images were collected automatically from web image searches and labelled with seven emotion categories, and the authors report human accuracy of roughly 65\%. We used the provided Training, PublicTest and PrivateTest splits without re-splitting. The Training split fitted all models, PublicTest was used for model comparison and tuning, and PrivateTest was used once per tuned model. All classifiers use seed 42, and the bootstrap intervals use 1,000 resamples (seed 0). Tasks 4 and 6 used the fuzzy environment (\texttt{environment-fuzzy.yml}: Python 3.10.22, NumPy 1.24.4, PyFUME 0.3.4, dlib 19.24.2, scikit-learn 1.7.2). Tasks 1, 2, 3 and 5 used the general environment (\texttt{environment-general.yml}: Python 3.12.15, scikit-learn 1.9.1, scikit-image 0.26.0 for HOG, dlib 20.0.1).

\textbf{Preprocessing and face detection:} FER2013 images are already cropped 48 $\times$ 48 grayscale faces, so no detection or resizing was applied. HOG was computed on the raw image, and PCA used pixels scaled to $[0, 1]$. For landmarks, the whole image is passed to \textit{dlib} as the face box, so there are no detection failures by construction. The cost is that landmark quality depends on the crop being well aligned, which fails for tilted, occluded or off-centre faces (Figure \ref{fig:err-gallery} in Appendix \ref{app:cm}).

\textbf{Quality checks:} Exact duplicates were found by MD5 hashing. 280 of 3,589 validation images (7.8\%) and 288 of 3,589 test images (8.0\%) also appear in the training split. We kept duplicates in all Task 1 evaluation splits, so results stay comparable to the provided protocol, and we report non-duplicated scores separately (Section \ref{sec:task1}). Duplicates were removed from the training set only in Task 5 (section \ref{sec:creative}).

\textbf{Limitations:} The dataset has strong class imbalance (Disgust is 1.5\% of training, See Figure \ref{fig:fer-dist} - Appendix \ref{app:cm}), noisy or ambiguous labels, low resolution, and cross-split leakage through duplicates. Figure \ref{fig:err-gallery} shows examples of the first two issues, with an infant labelled Angry and wide-eyed faces labelled Neutral.

\textbf{KDEF.} This database was presented by \textcite{goeleven_karolinska}, but for this assignment, we use only straight/frontal (S) faces, across seven comparable expressions (KDEF ``Afraid'' maps to FER ``fear''). Divide by identity: 80\% of subjects for training, 20\% for testing, with no subject overlap.

\textbf{Common representation:} Detect and crop KDEF faces with a margin preserving forehead and jaw; convert to grayscale, resize to $48\times48$ and apply the same normalization as FER. Apply identical feature extraction at training and inference. Fit any scaling, PCA or learned transformation only on training data.

\textbf{HOG features:} Histogram of Oriented Gradients (HOG) looks at the edges in an image and identifies where the brightness changes and in which direction. To extract the features, the $48\times48$ grayscale image is cut into a $6\times6$ grid of cells of $8\times8$ pixels. In each cell, the change in brightness is measured in 9 different directions ($0^\circ, 20^\circ, \dots, 160^\circ$). We then slide a $2\times2$ block across the grid, which results in $5\times5$ overlapping block positions:
\begin{equation*}
\text{HOG features}
= \underbrace{(5 \times 5)}_{\text{block positions}}
\times \underbrace{(2 \times 2)}_{\text{cells per block}}
\times \underbrace{9}_{\text{directions}}
= 900
\end{equation*}

\textbf{Fuzzy classifier protocol:}
For Task 4, geometric features were extracted directly from original-resolution frontal KDEF images, rather than the FER-compatible $48\times48$ images. Both acquisition sessions for each identity remained in the same partition. Model selection used training subjects only, preserving the held-out KDEF subjects for final evaluation.